\section{Divisors with small residues}\label{sec:residues}

We now construct the auxiliary denominator.
Its useful divisors make most residues small at each of a sequence of
scales.  At the smallest scales, several independent lists of divisors
ensure that every numerator has a suitable choice.

The underlying identity is elementary.  For positive integers \(M,Q,u,t\)
with \(t\mid M\), set
\[
 z=\left\lceil\frac{Qt}{u}\right\rceil,
 \qquad h=uz-Qt.
\]
Then \(z\ge1\), \(0\le h<u\), and
\begin{equation}\label{eq:residue-step}
 \frac{u}{MQ}=\frac{1}{(M/t)z}+\frac{h}{MQz}.
\end{equation}
Thus a small least nonnegative residue of \(-Qt\) modulo \(u\)
leaves a small numerator.  A zero residue finishes the expansion.

We use this identity with two denominator factors. At high levels we
take \(Q=C\): the large factor \(C\) supplies cancellation in
reciprocal phases as \(u\) varies. At lower levels we take \(Q=1\)
and construct expansions of \(u/M\). Dividing such an expansion by
\(C\) gives one of \(u/(MC)\), so the two regimes can be joined.
The switch below is made when the entire level cutoff is at most
\(e^S\); then \(u/M<1\) because our construction ensures \(M>e^S\).

\subsection{The denominator and the residue statement}

Fix the absolute constants
\begin{equation}\label{eq:residue-constants}
 \begin{gathered}
 K=100,\quad R=1000,\quad D_0=100000,\quad \eta=10^{-4},\\
 D_M=(2R+2)(K+2),\quad D_X=D_M+2,\quad D_C=4D_X.
 \end{gathered}
\end{equation}
Here \(K\) sets the prime scale, \(R\) is the number of independent
blocks, \(D_0\) sets the binary cutoff, and \(\eta\) sets the rate of
descent. These numerical choices leave room in the estimates below.
Let \(S\) tend to infinity, put \(m=\lfloor S/\log S\rfloor\), and fix
an integer
\begin{equation}\label{eq:C-range}
 e^{D_CS}\le C\le e^{2D_CS}.
\end{equation}
Define
\begin{equation}\label{eq:levels}
 \begin{gathered}
 \rho=e^{-\eta m},\qquad X_j=e^{D_XS}\rho^j,\\
 d=\min\{j\ge0:X_j\le e^m\},\qquad
 C_j=
 \begin{cases}
 C,&X_j>e^S,\\
 1,&X_j\le e^S.
 \end{cases}
 \end{gathered}
\end{equation}
These cutoffs need not be integers.  Every set or sum indexed by a
cutoff below consists of integers in the indicated range.
For \(0\le j\le d\), write
\begin{equation}\label{eq:least-residue}
 h_{j,t}(u)=u\left\lceil\frac{C_jt}{u}\right\rceil-C_jt.
\end{equation}
In particular, \(h_{j,t}(u)\) is the least nonnegative residue of
\(-C_jt\) modulo \(u\).

\begin{lemma}\label{lem:residues}
For all sufficiently large \(S\) and every integer \(C\) satisfying
\eqref{eq:C-range}, there are an integer \(M\) and indexed lists
\(T_0,T_1,\ldots,T_R\), each with \(2^m\) entries, such that:
\begin{enumerate}
\item
\(e^S<M\le e^{D_MS}\), and the least power of \(2\) at least
\(S^{D_0}\) divides \(M\).  Every entry of every \(T_i\) divides \(M\).
\item
For each \(0\le j<d\), let \(T=T_0\) if \(X_j>e^S\), and let
\(T=T_1\) otherwise.  Apart from at most
\(X_j e^{-c_*m}\) integers \(u\in(X_{j+1},X_j]\), at least
\(\rho |T|/2\) entries \(t\in T\) satisfy
\[
 h_{j,t}(u)\le X_{j+1},
 \qquad c_*=0.001.
\]
\item
For every integer \(S^{D_0}<u\le e^m\), some entry \(t\) of one of
\(T_1,\ldots,T_R\) satisfies
\[
 u\lceil t/u\rceil-t\le u^{1-\eta}.
\]
\end{enumerate}
Repeated values in a list are counted with their indices.
The threshold on \(S\) is independent of \(C\).
\end{lemma}

Here is the construction used to prove the lemma.
Let \(\mathcal P\) be the set of primes in \([S^K,2S^K]\), and put
\(P=|\mathcal P|\).  The prime number theorem gives
\(P\ge S^{K-1}\) for sufficiently large \(S\); see, for example,
\cite[p.~305, (1.1)]{Selberg1949}.
Choose distinct \(p_1,\ldots,p_m\in\mathcal P\),
and let
\[
 T_0=\left(\prod_{j=1}^m p_j^{I_j}\right)_{I\in\{0,1\}^m}.
\]
Its entries are distinct.  For each \(1\le i\le R\), choose
\(2m\) independent uniform samples
\(p_{i,j,\epsilon}\in\mathcal P\), where
\(1\le j\le m\) and \(\epsilon\in\{0,1\}\); all samples in different
blocks are also independent.  Set
\[
 \begin{gathered}
 T_i=\left(\prod_{j=1}^m p_{i,j,I_j}\right)_{I\in\{0,1\}^m},\\
 M=2^a\prod_{j=1}^m p_j
       \prod_{i=1}^R\prod_{j=1}^m p_{i,j,0}p_{i,j,1},
 \end{gathered}
\]
where \(2^a\) is the least power of \(2\) at least \(S^{D_0}\).
Sampling is with replacement.  Even when primes repeat, each list
entry uses a submultiset of the factors of \(M\), so it divides \(M\).

The size bounds hold for every realization:
\[
 \log M\ge Km\log S\ge K(S-\log S)>S,
\]
whereas
\[
 \log M
 \le D_0\log S+\log2+(2R+1)m(K\log S+\log2)
 \le D_MS
\]
eventually.  The last inequality follows because
\((2R+1)K<D_M\).  Every entry of every list also satisfies
\begin{equation}\label{eq:product-size}
 1\le t\le e^{(K+1)S}.
\end{equation}

We record the endpoint facts used below and in the descent.
For sufficiently large \(S\),
\begin{equation}\label{eq:level-endpoints}
 \begin{gathered}
 \frac{S}{2\log S}\le m\le\frac S{\log S},\\
 d=\left\lceil\frac{D_XS-m}{\eta m}\right\rceil
       \le K_d\log S,\qquad K_d=\frac{3D_X}{\eta},\\
 e^{(1-\eta)m}<X_d\le e^m,\qquad
 X_{j+1}\ge\sqrt{X_j}\quad(0\le j<d).
 \end{gathered}
\end{equation}
Indeed, for \(j<d\), \(\log X_j>m\), so
\(\log X_{j+1}=\log X_j-\eta m\ge(1-\eta)\log X_j\).
Moreover \(C_d=1\), and \(C_{j+1}\mid C_j\).

Table~\ref{tab:regimes} summarizes the roles of the lists. The upper
two rows are classified by the level cutoff \(X_j\), so a high level
may include some numerators below \(e^S\). The binary step is carried
out in Section~\ref{sec:descent}.
\begin{table}[ht]
\centering
\small
\begin{tabular}{@{}p{0.40\linewidth}p{0.28\linewidth}p{0.23\linewidth}@{}}
\toprule
Range and denominator factor & Divisors used & Coverage\\
\midrule
High levels: \(X_j>e^S\), \(C_j=C\)
 & Deterministic list \(T_0\) & Most numerators at each level\\[3pt]
Middle levels: \(e^m<X_j\le e^S\), \(C_j=1\)
 & The same random list \(T_1\) & Most numerators at each level\\[3pt]
Terminal range: \(S^{D_0}<u\le e^m\), factor \(1\)
 & Some list \(T_i\), \(1\le i\le R\) & Every numerator\\[3pt]
Binary range: \(1\le u\le S^{D_0}\), factor \(1\)
 & Powers of \(2\) dividing \(M\) & Every numerator\\
\bottomrule
\end{tabular}
\caption{The four stages of the numerator descent. The precise residue
guarantees are stated in Lemma~\ref{lem:residues}.}
\label{tab:regimes}
\end{table}

It remains to select the samples so that the two residue assertions
hold.  We first turn Fourier bounds into small residues, then prove
the needed bound for \(T_0\).
Lemma~\ref{lem:random} treats the random lists, and
Section~\ref{subsec:selection} makes the common choice.

\subsection{Discrepancy and the deterministic list}

We use the Erd\H{o}s--Tur\'an discrepancy inequality
\cite[Part~I, Theorem~III]{ErdosTuran1948}.
For any nonempty finite indexed list
\((x_t)_{t\in T}\) of real numbers, any interval \(J\subset[0,1)\),
and any integer \(H\ge1\), it states that
\begin{equation}\label{eq:erdos-turan}
 \left|
 \frac{\#\{t:\{x_t\}\in J\}}{|T|}-|J|
 \right|
 \ll \frac1H+
 \sum_{\ell=1}^H\frac1\ell
 \left|\frac1{|T|}\sum_{t\in T}\e(\ell x_t)\right|.
\end{equation}
The implied constant is absolute, and list multiplicities are retained.
The interval convention used below also retains entries at zero.
For a finite list, a sufficiently small common positive rotation sends
membership in $[0,\delta)$ to membership in $(0,\delta]$, including
all repeated endpoint entries correctly, while preserving every Fourier
modulus. Thus the interval convention in the original proof
\cite[Part~II, Sections~11--16]{ErdosTuran1948} gives the form needed here.

\begin{lemma}\label{lem:discrepancy}
For all sufficiently large real \(w\), suppose \(u,Q\) are positive
integers and \(T\) is a nonempty finite list of positive integers.
If
\[
 \left|\frac1{|T|}\sum_{t\in T}\e(\ell Qt/u)\right|
 \le e^{-3\eta w}
 \quad
 \left(1\le\ell\le\lfloor e^{4\eta w}\rfloor\right),
\]
then at least \(e^{-\eta w}|T|/2\) entries satisfy
\[
 0\le u\lceil Qt/u\rceil-Qt<e^{-\eta w}u.
\]
\end{lemma}
\begin{proof}
Apply \eqref{eq:erdos-turan} to \(x_t=-Qt/u\) and
\(J=[0,e^{-\eta w})\).  Conjugation leaves the Fourier bounds
unchanged.  The discrepancy is
\[
 O\bigl(e^{-4\eta w}+(1+4\eta w)e^{-3\eta w}\bigr)
 =o(e^{-\eta w}).
\]
It is therefore at most \(e^{-\eta w}/2\) eventually.
Finally,
\(\{-Qt/u\}=(u\lceil Qt/u\rceil-Qt)/u\), including when the
residue is zero.
\end{proof}

At a level \(X_j>e^S\), the large factor \(C\) makes the reciprocal
phases oscillate as \(u\) varies.  The following elementary estimate
will make that observation uniform.

\begin{lemma}\label{lem:reciprocal-phase}
Fix \(B\ge4\).  There are constants \(A_B,\delta_B>0\) such that,
for every sufficiently large \(U\), every real \(Z\) with
\(U^4\le |Z|\le U^B\), and every interval \(I\subset[U,2U]\),
\[
 \left|\sum_{n\in I}\e(Z/n)\right|\le A_BU^{1-\delta_B}.
\]
\end{lemma}
\begin{proof}
We give the derivative argument, including the elementary estimates
it uses.  This is the classical differencing method of van der Corput;
see \cite[Lemma~2.5]{GrahamKolesnik1991} for differencing and
\cite[Theorem~2.2]{GrahamKolesnik1991} for the second-derivative estimate.

First, if \(|a_n|\le1\) is supported on an interval of length at most
\(U\), then for \(2\le L\le U\),
\begin{equation}\label{eq:differencing}
 U^{-2}\left|\sum_n a_n\right|^2
 \ll L^{-1}
 +\max_{1\le h<L}U^{-1}
        \left|\sum_n a_{n+h}\overline{a_n}\right|.
\end{equation}
To see this, express the sum as
\(L^{-1}\sum_n\sum_{h=1}^L a_{n+h}\), apply Cauchy--Schwarz in \(n\),
and expand the square.  The outer support has length at most
\(U+L+2\), the diagonal contributes \(O(LU)\), and the off-diagonal
terms give \eqref{eq:differencing}.

We also need the second derivative bound
\begin{equation}\label{eq:second-derivative}
 \left|\sum_{n\in J}\e(g(n))\right|
 \ll_A U\sqrt\lambda+\lambda^{-1/2}+1
 \quad\text{if}\quad
 \lambda\le |g''(x)|\le A\lambda
\end{equation}
on an interval \(J\) of length at most \(U\), for \(g\in C^2(J)\).
Here is an elementary proof.  For \(\lambda\) above a fixed positive
constant the trivial bound suffices.  Otherwise put
\(\beta=\sqrt\lambda<1/4\).
The derivative \(g'\) is monotone and traverses a range of length
\(O_A(U\lambda)\).  The portions where \(g'\) is within \(\beta\)
of an integer have \(O_A(U\lambda+1)\) components and total length
\(O_A(U\beta+\beta/\lambda)\), since \(|g''|\ge\lambda\).
Their integer points contribute the same bound after adding the
number of components.

On each remaining interval \(g'\) is monotone and stays between
\(q+\beta\) and \(q+1-\beta\) for some integer \(q\).
The corresponding sum is \(O(\beta^{-1})\).  Indeed the increments
\(g(n+1)-g(n)\) are monotone in that same interval, and summation
by parts in
\[
 \e(g(n))=
 \frac{\e(g(n+1))-\e(g(n))}
      {\e(g(n+1)-g(n))-1}
\]
gives this bound: the reciprocal denominator has supremum and total
variation \(O(\beta^{-1})\).  Boundary terms add at most a constant
per interval.  Summing over the \(O_A(U\lambda+1)\) intervals proves
\eqref{eq:second-derivative}.

Now let \(k\) be a nearest integer to \(\log|Z|/\log U\), and set
\[
 Q_B=\lceil B\rceil+1,\qquad
 r=k-2,\qquad L=\lfloor U^{1/(10k)}\rfloor.
\]
There are only finitely many possible orders:
\begin{equation}\label{eq:derivative-order}
 4\le k\le Q_B,\qquad
 U^{-3/2}\le |Z|U^{-k-1}\le U^{-1/2}.
\end{equation}
Apply \eqref{eq:differencing} \(r\) times to
\(a_n=\mathbf1_I(n)\e(f(n))\), where \(f(x)=Z/x\), extending the
sequence by zero outside \(I\).  For positive shifts
\(h_1,\ldots,h_r<L\), the last phase is
\[
 g(x)=\Delta_{h_1}\cdots\Delta_{h_r}f(x),
 \qquad \Delta_hf(x)=f(x+h)-f(x).
\]
Its support is the interval on which both \(x\) and
\(x+h_1+\cdots+h_r\) belong to \(I\).
All intermediate points \(x+t_1+\cdots+t_r\), \(0\le t_i\le h_i\),
therefore remain in \([U,2U]\).  Empty supports contribute zero.

The fundamental theorem of calculus gives
\[
 g''(x)=
 \int_0^{h_1}\!\cdots\!\int_0^{h_r}
 f^{(k)}(x+t_1+\cdots+t_r)\,dt_r\cdots dt_1.
\]
Since \(f^{(k)}(x)=(-1)^k k!Zx^{-k-1}\) has constant sign,
\[
 \frac{k!}{2^{k+1}}\Lambda
 \le |g''(x)|\le k!\Lambda,
 \qquad \Lambda=|Z|U^{-k-1}\prod_{i=1}^r h_i.
\]
By \eqref{eq:derivative-order} and the choice of \(L\),
\[
 U^{-3/2}\le\Lambda
 \le U^{-1/2+(k-2)/(10k)}\le U^{-2/5}.
\]
Thus \eqref{eq:second-derivative} bounds each last correlation by
\(O_k(U^{4/5}+U^{3/4}+1)=O_k(U^{4/5})\).

For clarity, normalize every correlation by \(U\), and let \(\sigma_j\)
be their maximum after \(j\) shifts.  Then
\[
 \sigma_{k-2}\ll_k U^{-1/5},
 \qquad \sigma_j^2\ll L^{-1}+\sigma_{j+1}.
\]
Induction, using \(L\ge U^{1/(10k)}/2\) for large \(U\), yields
\[
 \sigma_0\ll_k U^{-1/(10k\,2^{k-2})}.
\]
Taking the largest implied constant over \(4\le k\le Q_B\) and,
for example, \(\delta_B=(10Q_B2^{Q_B})^{-1}\) proves the lemma.
In particular, although \(k\) depends on \(Z,U\), the final constants
depend only on \(B\).
\end{proof}

\begin{lemma}\label{lem:deterministic}
For the deterministic list \(T_0\) constructed above, every level
\(0\le j<d\) with \(X=X_j>e^S\), and every integer
\(1\le\ell\le H:=\lfloor e^{4\eta m}\rfloor\), one has
\begin{equation}\label{eq:deterministic-mean}
 \frac1X\sum_{X_{j+1}<u\le X}
 \left|\frac1{|T_0|}\sum_{t\in T_0}\e(\ell Ct/u)\right|^2
 \le e^{-0.01m}
\end{equation}
for all sufficiently large \(S\), uniformly in \(C\).
\end{lemma}
\begin{proof}
Write \(Y=X_{j+1}=\rho X\); then
\(Y>e^{S-\eta m}\ge e^{0.9S}\).
Expanding the square, the diagonal contributes at most \(2^{-m}\)
because the \(2^m\) subset products are distinct.
For an off-diagonal pair put \(Z=\ell C(t-t')\).
The difference \(t-t'\) is a nonzero integer, and
\eqref{eq:C-range} and \eqref{eq:product-size} give
\[
 e^{D_CS}\le |Z|\le e^{(2D_C+K+2)S}.
\]
Split \((Y,X]\) into intersections with dyadic intervals \([U,2U]\).
We may take \(Y/2\le U\le X\), so for large \(S\),
\[
 e^{0.8S}\le U\le e^{D_XS},\qquad
 U^4\le |Z|\le U^{3D_C}.
\]
Lemma~\ref{lem:reciprocal-phase}, with the fixed value \(B=3D_C\),
bounds the sum over each piece by
\(O(Ue^{-0.8\delta_BS})\).
The sum of the dyadic \(U\)'s is \(O(X)\).  Thus
\[
 \left|\sum_{Y<u\le X}\e(Z/u)\right|\ll Xe^{-cS}
\]
for an absolute \(c>0\).  Averaging over all off-diagonal pairs
cancels their number against \(|T_0|^2\).  The left side of
\eqref{eq:deterministic-mean} is consequently at most
\[
 2^{-m}+O(e^{-cS})\le e^{-0.01m}
\]
eventually, since \(S/m\to\infty\).
\end{proof}

The deterministic estimate controls the high levels for every
permitted \(C\).  We next prove the random estimate for the lower
levels and for individual terminal numerators.

\subsection{Random products}\label{subsec:random}

For the smaller moduli, two independent products of sampled primes supply
the cancellation.  Their integer values are smaller than a suitable
divisor of the modulus, so unique factorization controls their residue
distributions.

\begin{lemma}\label{lem:random}
Let \(K=100\), \(D_0=100000\), and \(m=\lfloor S/\log S\rfloor\).
Let \(\mathcal P\) be a set of primes in \([S^K,2S^K]\) with
\(P=|\mathcal P|\ge S^{K-1}\), and sample
\[
 p_{j,\epsilon}\qquad
 (1\le j\le m,\ \epsilon\in\{0,1\})
\]
independently and uniformly from \(\mathcal P\).  For
\(I\in\{0,1\}^m\), put \(t_I=\prod_{j=1}^m p_{j,I_j}\).
For every integer \(u\) such that
\[
 D_0\log S\le V:=\log u\le S,
 \qquad w=\min(m,V),
\]
and every integer \(1\le l\le \exp(0.005w)\), one has
\begin{equation}\label{eq:random-second-moment}
 \E\left|2^{-m}\sum_{I\in\{0,1\}^m}\e(lt_I/u)\right|^2
 \le \exp(-0.01w)
\end{equation}
for all sufficiently large \(S\).  The onset is absolute and uniform in
\(u,l,\mathcal P\).
\end{lemma}

\begin{proof}
Expanding the square gives
\[
 4^{-m}\sum_{I,J\in\{0,1\}^m}
       \E\,\e\bigl(l(t_I-t_J)/u\bigr).
\]
For uniformly counted pairs \(I,J\), their Hamming distance \(B\) has
distribution \(\operatorname{Bin}(m,1/2)\).  The exponential Markov
inequality gives
\[
 \Prob(B<m/4)
 \le 2^{m/4}\E\,2^{-B}
 =2^{m/4}(3/4)^m
 \le e^{-m/16}.
\]
We bound these pairs trivially.  Fix any remaining pair, and set
\[
 s=\left\lfloor\frac{0.75V}{K\log S}\right\rfloor.
\]
Since \(V/(K\log S)\ge1000\), and \(V\le S\), we have
\begin{equation}\label{eq:random-product-lengths}
 \frac{0.74V}{K\log S}\le s\le\frac{0.75V}{K\log S},
 \qquad 2s\le m/4
\end{equation}
for sufficiently large \(S\).

Select the first \(2s\) differing coordinates in increasing order, and
divide them into two sets of \(s\) coordinates each.  This choice depends
only on \(I,J\), not on the sampled primes.  Expose every labelled prime
sample except the \(I\)-selected sample at these \(2s\) coordinates.
Then \(t_J\) is fixed and
\[
 t_I=U_0U_1U_2,
\]
where \(U_0\) is fixed by the exposure and \(U_1,U_2\) are independent
products of \(s\) fresh uniform samples each.  Equal prime values cause
no difficulty: the \(I\)-selected and \(J\)-selected variables at a
differing coordinate are distinct independent samples.

Write
\[
 g=\gcd(lU_0,u),\qquad q=u/g.
\]
The coefficient \(lU_0/u\) reduces to a fraction with denominator
\(q\). We will show that \(q\) is usually large enough for distinct
integer values of each fresh product to remain distinct modulo \(q\).
Their small point probabilities will then give cancellation by
additive-character orthogonality.

We first bound the probability, over the exposed variables, that
\(q<u^{0.9}\).  Since \(l\le u^{0.005}\) and
\(\gcd(lU_0,u)\le l\gcd(U_0,u)\), this event implies
\(\gcd(U_0,u)>u^{0.095}\).
Call a sample contributing to \(U_0\) a hit if its prime value divides
\(u\), and count hits with their sample multiplicities.  If their number
is \(B_0\), then
\[
 \gcd(U_0,u)\le \prod_{\text{hit samples }p}p
       \le S^{(K+1)B_0}.
\]
Thus failure requires at least
\[
 k_0=\left\lceil\frac{0.08V}{(K+1)\log S}\right\rceil
\]
hits.  At most \(S\) primes in the sampling interval divide \(u\), so
each sample has hit probability at most
\(S/P\le S^{-(K-2)}\).  There are at most \(m\le S\) independent samples
in \(U_0\).  A union bound over sets of \(k_0\) positions yields
\begin{equation}\label{eq:random-gcd}
 \Prob(q<u^{0.9})
 \le S^{-(K-3)k_0}
 \le \exp\left(-\frac{0.08(K-3)}{K+1}V\right)
 \le u^{-0.05}.
\end{equation}
If \(k_0>m-2s\), the event is empty.  This argument allows prime powers
in \(u\) and repeated sampled primes.

Now fix an exposure with \(q\ge u^{0.9}\).  Every possible value of
either fresh product is at most
\[
 (2S^K)^s\le S^{(K+1)s}\le u^{0.7575}<u^{0.8}<q.
\]
By unique factorization, an integer can arise as the product of at most
\(s!\) ordered \(s\)-tuples of primes.  Consequently each atom in the
distribution of \(U_i\), for \(i=1,2\), has probability at most
\[
 \frac{s!}{P^s}\le S^{-(K-2)s}
      \le u^{-0.7252}\le u^{-0.7}.
\]
Reduction modulo \(q\) is injective on these integer values.  The
probability vectors \(\alpha,\beta\) of \(U_1,U_2\) on
\(\Z/q\Z\) therefore satisfy
\[
 \|\alpha\|_2,\|\beta\|_2\le u^{-0.35}.
\]

The reduced fraction \(lU_0/u=c/q\) has \(\gcd(c,q)=1\).
Complete additive-character orthogonality gives
\[
 \sum_{x\bmod q}\left|\sum_{y\bmod q}\beta_y\e(cxy/q)\right|^2
       =q\sum_{y\bmod q}|\beta_y|^2.
\]
Indeed the inner sum over \(x\) vanishes unless
\(q\mid c(y-y')\), equivalently \(y=y'\) modulo \(q\).
This identity holds for composite \(q\) as well.  Cauchy--Schwarz gives
\[
 \left|\sum_{x,y\bmod q}\alpha_x\beta_y\e(cxy/q)\right|
 \le \sqrt q\,\|\alpha\|_2\|\beta\|_2
 \le u^{-0.2}.
\]
The factor \(\e(-lt_J/u)\) is fixed under the exposure.  Averaging this
conditional estimate and using \eqref{eq:random-gcd} bounds the
contribution of our fixed pair \(I,J\) in absolute value by
\(u^{-0.05}+u^{-0.2}\).

Combining this with the Hamming exception, we obtain
\[
 \E\left|2^{-m}\sum_I\e(lt_I/u)\right|^2
 \le e^{-m/16}+u^{-0.05}+u^{-0.2}
 \le 3e^{-0.05w}
 \le e^{-0.01w}.
\]
The final inequality is uniform because
\(\min(m,D_0\log S)\) tends to infinity with \(S\).
\end{proof}

\subsection{A simultaneous choice of divisors}\label{subsec:selection}

We now choose all random blocks at once.  One block will suffice at all
intermediate levels outside small exceptional sets; the independent
blocks together will cover every terminal integer.

\begin{proof}[Proof of Lemma~\ref{lem:residues}]
The construction already gives the required size of \(M\), its power of
two, and divisibility of every list entry, for every realization of the
samples.  We prove the two residue properties.

Put \(H=\lfloor e^{4\eta m}\rfloor\).
At a level with \(X_j>e^S\), Lemma~\ref{lem:deterministic} and Markov's
inequality show that the number of integers
\(u\in(X_{j+1},X_j]\) for which
\[
 \left|2^{-m}\sum_{t\in T_0}\e(lCt/u)\right|>e^{-3\eta m}
\]
at some \(1\le l\le H\) is at most
\[
 X_j H e^{6\eta m}e^{-0.01m}
 \le X_j e^{-0.009m}.
\]
Conjugation changes the Fourier sign without changing the absolute
value.  Lemma~\ref{lem:discrepancy}, with scale \(m\), therefore gives at
least \(\rho|T_0|/2\) entries with
\[
 h_{j,t}(u)<\rho u\le X_{j+1}
\]
outside that exceptional set.  This holds separately at every such
level and requires no probabilistic selection.

Next suppose \(j<d\) and \(X_j\le e^S\).  Then \(X_j>e^m\), and every
\(u\in(X_{j+1},X_j]\) satisfies
\[
 (1-\eta)m<\log u\le S,\qquad
 w=\min(m,\log u)\ge(1-\eta)m.
\]
For sufficiently large \(S\), these integers satisfy the lower bound
in Lemma~\ref{lem:random}.  Moreover
\(4\eta m\le0.005w\), so that Lemma applies to all \(1\le l\le H\).
In block \(1\), Markov's inequality followed by the frequency union
bound shows that
\begin{align*}
 &\Prob\left(
   \max_{1\le l\le H}
   \left|2^{-m}\sum_{t\in T_1}\e(lt/u)\right|>e^{-3\eta m}
 \right)\\
 &\hspace{2em}\le
 \exp\bigl((10\eta-0.01(1-\eta))m\bigr)
 =e^{-0.008999m}\le e^{-0.005m}.
\end{align*}
Let \(B_j\) count integers in this level that fail this Fourier bound.
Linearity of expectation and another application of Markov give
\[
 \E B_j\le X_j e^{-0.005m},
 \qquad
 \Prob(B_j>X_j e^{-0.001m})\le e^{-0.004m}.
\]
As \(d=O(\log S)\), with probability \(1-o(1)\) these bounds hold at
every level under consideration.  No independence between different
integers or levels is needed.  Here \(C_j=1\), so
Lemma~\ref{lem:discrepancy} again gives the required
\(\rho|T_1|/2\) entries outside the exceptional sets.

Finally fix an integer \(S^{D_0}<u\le e^m\).
Now \(w=\log u\), and Lemma~\ref{lem:random}, Markov's inequality and
the union over \(1\le l\le\lfloor u^{4\eta}\rfloor\) give
\begin{align*}
 &
 \Prob\left(
  \max_{1\le l\le\lfloor u^{4\eta}\rfloor}
   \left|2^{-m}\sum_{t\in T_i}\e(lt/u)\right|>u^{-3\eta}
 \right)\\
 &\hspace{2em}\le u^{10\eta-0.01}
 =u^{-0.009}\le u^{-0.005}
\end{align*}
in each block \(i\).  On success, Lemma~\ref{lem:discrepancy}, with
scale \(\log u\), supplies a list entry for which
\[
 u\lceil t/u\rceil-t<u^{1-\eta}.
\]
The \(R=1000\) blocks are independent for this fixed \(u\).
The probability that every block fails is therefore at most \(u^{-5}\).
A union bound over all terminal integers shows that
\[
 \Prob(\text{some terminal integer has no successful block})
 \le \sum_{u>S^{D_0}}u^{-5}=o(1).
\]
Thus every terminal integer has a suitable divisor, simultaneously,
with probability tending to one.

The sum of this failure probability and the middle-level failure
probability is \(o(1)\).  Hence a realization satisfying both sets of
requirements exists.  Fix it and the resulting \(M\).
Since \(e^{-0.009m}\le e^{-0.001m}\), the same
\(c_*=0.001\) is valid at every level.  This proves both properties for
one common \(M\), for the fixed \(S\) and \(C\).
\end{proof}

